\chapter{Mood disorders}
\index{Mood disorders}

\section*{Definition and Overview}
Mood disorders, also referred to as affective disorders, are a category of mental health conditions characterised by significant, persistent disturbances in a person's emotional state or mood. These disorders primarily include major depressive disorder (unipolar depression), bipolar I and II disorders (distinguished by manic or hypomanic episodes interspersed with depressive phases), persistent depressive disorder (dysthymia), and cyclothymic disorder. Unlike transient emotional fluctuations in response to daily life events, mood disorders involve sustained changes in affect that are disproportionate to circumstances and lead to profound functional impairment, distress, and alterations in sleep, energy, cognition, and behaviour.

\section*{Epidemiology}
Mood disorders are among the most common psychiatric conditions globally and represent a leading cause of disability worldwide. In many countries, national surveys indicate that approximately 15--20\% of the population will experience a mood disorder at some point in their lifetime. Major depressive disorder is twice as common in females as in males, whereas bipolar disorder has a relatively equal sex distribution, affecting around 1--2\% of the population. The peak age of onset for most mood disorders is in late adolescence or early adulthood (typically between 15 and 30 years), though they can manifest at any stage of life. High rates of comorbidity are observed with anxiety disorders, substance use disorders, and chronic physical health conditions.

\section*{Aetiology and Pathophysiology}
The development of mood disorders is multifactorial, arising from a complex interplay of genetic, biological, environmental, and psychological factors.
\begin{itemize}
    \item \textbf{Genetic factors:} Family and twin studies show high heritability rates, particularly for bipolar disorder (estimated heritability of 70--80\%) and to a lesser extent major depressive disorder (35--40\%).
    \item \textbf{Neurochemical dysregulation:} Classical monoamine hypothesis suggests altered transmission and receptor sensitivity of neurotransmitters, primarily serotonin, noradrenaline, and dopamine.
    \item \textbf{Neuroendocrine dysfunction:} Hyperactivity of the hypothalamic-pituitary-adrenal (HPA) axis, leading to elevated cortisol levels and impaired feedback mechanisms, which is commonly observed in severe depression.
    \item \textbf{Neuroplasticity and brain structure:} Reduced levels of brain-derived neurotrophic factor (BDNF) and structural alterations (such as volume reduction) in key brain regions like the prefrontal cortex, hippocampus, and amygdala.
    \item \textbf{Psychosocial stressors:} Adverse childhood experiences, chronic life stress, trauma, grief, and lack of social support can trigger or exacerbate episodes in vulnerable individuals.
\end{itemize}

\section*{Clinical Features}
The presentation depends on the specific disorder and whether the patient is experiencing a depressive, manic, or hypomanic episode.
\begin{itemize}
    \item \textbf{Depressive symptoms:} Persistent low mood, feelings of sadness or emptiness, anhedonia (loss of interest or pleasure in all activities), significant weight changes, insomnia or hypersomnia, psychomotor agitation or retardation, fatigue, feelings of worthlessness, impaired concentration, and suicidal ideation.
    \item \textbf{Manic symptoms:} Persistently elevated, expansive, or irritable mood, inflated self-esteem or grandiosity, decreased need for sleep (e.g., feeling rested after 3 hours), pressured speech, flight of ideas, distractibility, increased goal-directed activity, and reckless or impulsive behaviour (e.g., excessive spending or unsafe driving).
    \item \textbf{Hypomanic symptoms:} Similar to manic symptoms but less severe, lasting at least four consecutive days, and not causing marked impairment in social or occupational functioning or requiring hospitalisation.
    \item \textbf{Somatic complaints:} Chronic pain, headaches, gastrointestinal disturbances, and fatigue without an identifiable physical cause.
\end{itemize}

\section*{Differential Diagnosis}
It is critical to distinguish primary mood disorders from other medical conditions and psychiatric disorders that present with mood fluctuations.
\begin{itemize}
    \item \textbf{Endocrine disorders:} Hypothyroidism (mimics depression) or hyperthyroidism (mimics mania or anxiety), which must be ruled out via laboratory testing.
    \item \textbf{Substance-induced mood disorder:} Mood changes directly caused by the physiological effects of alcohol, illicit drugs (e.g., amphetamines causing mania), or medications (e.g., corticosteroid-induced depression or psychosis).
    \item \textbf{Neurological conditions:} Early-stage dementia, Parkinson's disease, multiple sclerosis, or stroke, which can present with apathy, depression, or emotional lability.
    \item \textbf{Schizoaffective disorder:} Characterised by prominent psychotic symptoms (hallucinations or delusions) occurring in the absence of a mood episode, whereas in mood disorders, psychotic features only occur within a depressive or manic episode.
    \item \textbf{Borderline personality disorder:} Features rapid, transient mood shifts in response to interpersonal stressors, differing from the sustained mood episodes of bipolar or unipolar depression.
\end{itemize}

\section*{When to Seek Medical Care}
Individuals should seek medical or psychological evaluation if they experience persistent changes in mood that interfere with their daily activities, relationships, or work.

\textbf{Contact your local emergency services for:}
\begin{itemize}
    \item Active suicidal ideation with a plan or intent, or self-harm behaviour.
    \item Severe manic states where the individual is unable to care for themselves or represents a danger to themselves or others.
    \item Acute psychosis (delusions or hallucinations) accompanying a severe mood episode.
\end{itemize}

\section*{Diagnosis and Investigations}
Diagnosis is based on a detailed clinical interview, psychiatric history, and standardised diagnostic criteria (such as DSM-5-TR or ICD-11).
\begin{itemize}
    \item \textbf{Psychiatric history and mental state examination (MSE):} Assessment of mood, affect, speech, thought content, perception, cognition, and risk.
    \item \textbf{Diagnostic criteria checklist:} Evaluation of the duration, number, and severity of symptoms to distinguish between Major Depressive Disorder, Bipolar I/II, and dysthymia.
    \item \textbf{Laboratory screening:} Full blood count, thyroid function tests (TFTs), liver function tests, electrolytes, vitamin B12, and vitamin D levels to rule out organic causes.
    \item \textbf{Urine toxicology screen:} Indicated to identify substance use that may be causing or complicating the clinical presentation.
    \item \textbf{Collateral history:} Interviewing family members or close contacts, particularly important for identifying past hypomanic or manic episodes that the patient may not recognise.
\end{itemize}

\section*{Management}
A multimodal, collaborative treatment plan incorporating pharmacotherapy, psychotherapy, and lifestyle modifications is standard.
\begin{itemize}
    \item \textbf{Psychotherapy:} Cognitive Behavioural Therapy (CBT), Acceptance and Commitment Therapy (ACT), or Interpersonal Therapy (IPT) for depressive disorders; Family-Focused Therapy (FFT) or Interpersonal and Social Rhythm Therapy (IPSRT) for bipolar disorders.
    \item \textbf{Pharmacotherapy:} Antidepressants (SSRIs, SNRIs) for depression; mood stabilisers (lithium, sodium valproate) and atypical antipsychotics (quetiapine, olanzapine) for bipolar disorder.
    \item \textbf{Electroconvulsive therapy (ECT):} Highly effective neuromodulation therapy reserved for severe, treatment-resistant depression, severe mania, or catatonia.
    \item \textbf{Lifestyle interventions:} Structuring regular sleep patterns (crucial in bipolar disorder), engaging in regular physical exercise, avoiding alcohol and recreational drugs, and establishing a support system.
    \item \textbf{Psychoeducation:} Educating the patient and family about symptom monitoring, medication adherence, and early warning signs of relapse.
\end{itemize}

\section*{Complications}
Mood disorders carry serious risks of long-term physical, mental, and social consequences.
\begin{itemize}
    \item \textbf{Suicide and self-harm:} The most critical complication, with the lifetime risk of suicide being significantly elevated in both severe depression and bipolar disorder.
    \item \textbf{Substance abuse:} Self-medication with alcohol or drugs, leading to secondary substance use disorders and poorer prognosis.
    \item \textbf{Cardiovascular and metabolic disease:} Increased risk of obesity, diabetes, and coronary artery disease, mediated by HPA-axis hyperactivity, lifestyle factors, and medication side effects.
    \item \textbf{Socioeconomic impairment:} Chronic unemployment, financial difficulties, homelessness, and breakdown of personal relationships.
    \item \textbf{Cognitive decline:} Persistent deficits in executive function, memory, and attention, particularly following recurrent, untreated manic or depressive episodes.
\end{itemize}

\section*{Prognosis and Outcomes}
The prognosis varies depending on the age of onset, specific diagnosis, severity, and adherence to treatment. Major depressive disorder is often recurrent, with approximately 50--80\% of individuals experiencing a second episode. Bipolar disorder is typically a lifelong, chronic cyclic condition requiring long-term maintenance treatment. However, with a combination of appropriate medication, regular psychotherapy, and lifestyle management, the vast majority of patients achieve symptom remission, regain social and occupational function, and lead fulfilling lives. Early intervention and a strong therapeutic alliance are key predictors of favourable long-term outcomes.

\section*{Prevention and Self-Care}
Self-management and prevention focus on early detection of symptoms and maintaining stability.
\begin{itemize}
    \item \textbf{Maintain a sleep routine:} Go to bed and wake up at the same time daily, as sleep disruption is a major trigger for manic and depressive episodes.
    \item \textbf{Develop a relapse signature plan:} Document early warning signs of mood changes (e.g., changes in sleep, irritability) and have a pre-planned action checklist.
    \item \textbf{Avoid mood-altering substances:} Limit alcohol and caffeine intake, and completely avoid illicit stimulants or cannabis.
    \item \textbf{Engage in regular physical activity:} Aim for at least 30 minutes of moderate exercise most days to promote neuroplasticity and improve mood.
    \item \textbf{Stay connected:} Maintain a strong social network and attend support groups to prevent isolation during low mood phases.
\end{itemize}

\section*{Sources and Further Reading}
The following reputable organisations provide further information and support for mood disorders:
\begin{itemize}
    \item Beyond Blue. \textit{Depression and bipolar disorder.} https://www.beyondblue.org.au/
    \item Black Dog Institute. \textit{Mood disorders and treatments.} https://www.blackdoginstitute.org.au/
    \item Headspace. \textit{Understanding depression for young people.} https://headspace.org.au/
    \item Mayo Clinic. \textit{Depression and bipolar support.} https://www.mayoclinic.org/
    \item Royal Australian and New Zealand College of Psychiatrists (RANZCP). \textit{Clinical practice guidelines for mood disorders.} https://www.ranzcp.org/
\end{itemize}
